\section{Quantitative first-energy convergence in every local capacity}
\label{sec:r11-first-energy-capacity-topology}

Retain the strict analytic box, background extension and
$\mathcal E^1$ hypotheses of
Section~\ref{sec:r11-global-first-energy}.
All functions below are actual negative classical members
of that class; all capacities use the full original tester
class. Put $d=q+1$ and $\alpha_q=2^{-q}$.
The symmetric discrepancy $\mathscr I_S$ is (GE18).
At a joint pole, a difference is used only as a Borel
coefficient representative; all involved products and
capacities annihilate that pole union.

The bounded common-collar predecessors are the frozen R3
\texttt{energy\_convex\_blocks.tex},
\texttt{thm:r3-I-capacity}, and
\texttt{strong\_relative\_topology.tex}.
They already contain the positive discrepancy, the mixed
energy transfer recurrence, and bounded relative energy
topologies. The new scope here is the global unbounded
first-energy class, uniform control over its energy testers,
the analytic all-local domain extension, and critical
endpoint continuity without a prescribed common collar.

\begin{theorem}[A uniform first-energy tester estimate]
\label{thm:r11-first-energy-uniform-tester}
For $u,v,\psi\in\mathcal E^1$, put $w=u+v$ and
$M=\max\{e_1(w),e_1(\psi)\}$. Then
\[
 \boxed{\quad
 \int_V|u-v|\,d\mu_\psi
   \leq 2^{(q-1)\alpha_q}
          \mathscr I_S(u,v)^{\alpha_q}M^{1-\alpha_q}.
 \quad}
 \tag{EV1}
\]
The conclusion is zero if $M=0$. For a mixed top measure
in place of $\mu_\psi$, the positive expansion of the sum
of its $q$ potentials gives the same estimate divided by
$q!$, with that sum as $\psi$.
\end{theorem}

\noindent\textbf{Status: \texttt{NEW\_PROVED}} in the
unbounded global first-energy class.

\begin{proof}
Write $h=u-v$. For $0\leq k\leq q-1$ set
\[
 J_k=\int_V\Gamma_{N(\psi^{[k]},w^{[q-1-k]})}(h).
 \tag{EV2}
\]
These are finite nonnegative numbers by (GE7). The positive
expansion of $Dw=Du+Dv$ and (GE18) give
$J_0\leq2^{q-1}\mathscr I_S(u,v)$.
For $k\leq q-2$, full mixed energy symmetry exchanges one
$\psi$ slot with one $h$ slot. Expanding that latter
$h=u-v$ in the background gives
\[
 \begin{split}
 J_{k+1}
 &=\int\Gamma_{N(u,\psi^{[k]},w^{[q-2-k]})}(h,\psi)
   -\int\Gamma_{N(v,\psi^{[k]},w^{[q-2-k]})}(h,\psi)\\
 &\leq\left[J_k\,
       \mathcal I_S(\psi^{[k+2]},w^{[q-1-k]})\right]^{1/2}
 \leq(J_kM)^{1/2}.
 \end{split}
 \tag{EV3}
\]
For the middle inequality apply energy Cauchy--Schwarz in
each of the two positive backgrounds and then Cauchy--Schwarz
to the resulting two nonnegative summands. Their first
diagonal masses add to $J_k$, and their second diagonal
masses add to the displayed mixed energy. The last inequality
is the constant-one Holder inequality (GE10). Thus
\[
 J_{q-1}\leq J_0^{2^{-(q-1)}}M^{1-2^{-(q-1)}}.
 \tag{EV4}
\]
When $q=1$ this is just the identity $J_{q-1}=J_0$.

The real-span element
$|h|=2\max(u,v)-u-v$ has the exact energy identity
$\Gamma_Q(|h|)=\Gamma_Q(h)$ for every background $Q$
in (EV2). To justify it in the present class, use canonical
bounded truncations of $u,v,\max(u,v)$ and the background
slots. At level $j$ one has
$\max(u_j,v_j)=\max(\max(u,v),-j)$.
The bounded absolute-value chain and the zero energy of
the level $\{u_j=v_j\}$ prove the identity at that stage.
Endpoint local total-variation convergence (PG11) passes
it to the unbounded currents. Both are already globally
finite by (GE7).

Global Green, followed by Cauchy--Schwarz, now yields
\[
 \int|h|\,d\mu_\psi
  =-\int\Gamma_{N(\psi^{[q-1]})}(|h|,\psi)
       \leq J_{q-1}^{1/2}e_1(\psi)^{1/2}.
 \tag{EV5}
\]
Combine this with (EV4), the bound on $J_0$, and
$e_1(\psi)\leq M$. If $M=0$, the mixed Holder estimate
already makes the integral zero. The mixed-measure
assertion follows from
$q!M(\psi_1,\ldots,\psi_q)\leq\mu_{\sum_i\psi_i}$.
\end{proof}



\begin{theorem}[Quantitative capacity control on every compact/domain pair]
\label{thm:r11-first-energy-all-local-capacity}
Fix a compact $K\Subset V$ and choose $A_K>0$ so that
$A_K\rho<-2$ on a neighborhood of $K$. Set
\[
 B_K=2A_K^q\mu_\rho(V),\qquad
 \Lambda_K(u,v)=\max\{e_1(u+v),B_K\}.
 \tag{EV6}
\]
For $\delta>0$ put
$G_\delta(u,v)=\{u<v-\delta\}\cup\{v<u-\delta\}$,
using actual extended comparisons. Then
\[
 C_S(K\cap G_\delta(u,v);V)
    \leq\frac{2^{(q-1)\alpha_q}}\delta
          \Lambda_K(u,v)^{1-\alpha_q}
                      \mathscr I_S(u,v)^{\alpha_q}.
 \tag{EV7}
\]
If $K\subset\operatorname{Int}_X K'\Subset V$, the same
bound for $K'$ bounds the outer capacity of
$K\cap G_\delta(u,v)$.

Consequently for every ordinary $\Omega\subset V$ and
compact $K\Subset\Omega$, and every $M_0<\infty$, there
is a finite constant $C_{K,\Omega,M_0,S}$ such that
$e_1(u),e_1(v)\leq M_0$ implies
\[
 C_{S|\Omega}^*(K\cap G_\delta(u,v);\Omega)
       \leq\frac{C_{K,\Omega,M_0,S}}\delta
                            \mathscr I_S(u,v)^{\alpha_q}.
 \tag{EV8}
\]
This is the complete all-local capacity topology, by the
proved analytic affinoid domain comparison, not an assumed
identification of two admissibility domains.
\end{theorem}

\noindent\textbf{Status: \texttt{NEW\_PROVED}} as a
global-energy estimate for all original local testers.

\begin{proof}
For a normalized tester $0\leq a\leq1$ on $V$ put
$\psi=\max(a-2,A_K\rho)$.
It is bounded in $[-2,0]$, has compact negative sublevels,
and agrees with $a-2$ near $K$.
Comparison with $cA_K\rho$ for $c>1$ gives
$\mu_\psi(V)\leq A_K^q\mu_\rho(V)$ after $c\downarrow1$.
Thus $\psi\in\mathcal E^1$ and $e_1(\psi)\leq B_K$,
uniformly in the original tester $a$. Ordinary locality
gives $\mathbf1_K\mu_a=\mathbf1_K\mu_\psi$.
The joint pole set has raw capacity zero, so it has zero
mass for $\mu_a$. On its complement, Chebyshev and (EV1)
give
\[
 \mu_a(K\cap G_\delta(u,v))
       \leq\delta^{-1}\int|u-v|\,d\mu_\psi
       \leq\frac{2^{(q-1)\alpha_q}}\delta
             \Lambda_K(u,v)^{1-\alpha_q}
                          \mathscr I_S(u,v)^{\alpha_q}.
\]
Take the full tester supremum to prove (EV7).

To obtain the outer bound, fix $\eta>0$. Outer
quasicontinuity on $K'$ and outer nullity of the poles
give an ordinary open $O$ with $C_S(O;V)<\eta$ such
that both potentials are finite and continuous on
$K'\setminus O$. In the relative space
$\operatorname{Int}_X K'\setminus O$ the strict gap is
open. Hence there is an ordinary open $H\subset
\operatorname{Int}_X K'$ such that
$H\setminus O=G_\delta(u,v)\cap
(\operatorname{Int}_X K'\setminus O)$.
The open set $H\cup O$ covers $K\cap G_\delta$ and is
contained in $(K'\cap G_\delta)\cup O$. Therefore
\[
 C_S^*(K\cap G_\delta;V)
              \leq C_S(K'\cap G_\delta;V)+\eta.
\]
Let $\eta\downarrow0$ and use (EV7) on $K'$.

Finally choose $K'$ inside $\Omega$ and an ordinary open
core containing $K$ compactly in $\Omega$. The outer
domain comparison (AL9) multiplies this estimate by a fixed
geometric constant. Minkowski (GE11) gives
$e_1(u+v)\leq2^{q+1}M_0$, proving (EV8).
No general Borel raw/outer equality is used.
\end{proof}

\begin{theorem}[The strong first-energy topology]
\label{thm:r11-strong-first-energy-topology}
For any $u,v\in\mathcal E^1$,
\[
 \left|e_1(u)^{1/d}-e_1(v)^{1/d}\right|^d
                                      \leq\mathscr I_S(u,v).
 \tag{EV9}
\]
For a sequence $u_j\in\mathcal E^1$ and an actual
$u\in\mathcal E^1$, the following are equivalent:
\begin{enumerate}
\item $\mathscr I_S(u_j,u)\to0$.
\item $u_j\to u$ in all-local truncated original capacity
and $e_1(u_j)\to e_1(u)$.
\end{enumerate}
The forward implication gives outer-capacity convergence
of the full actual gaps on every compact/domain pair.
This theorem describes convergence to an already specified
actual class member; it does not assert classical
representability of every Cauchy sequence.
\end{theorem}

\noindent\textbf{Status: \texttt{NEW\_PROVED}} for the
unbounded class. The bounded fixed-collar energy topology
was already frozen in R3.

\begin{proof}
Put $a=e_1(u)^{1/d}$ and $b=e_1(v)^{1/d}$.
Expand the definition (GE18), then use the constant-one
mixed Holder inequality on its two cross terms:
\[
 \begin{split}
 \mathscr I_S(u,v)
 &=e_1(u)+e_1(v)
       -\mathcal I_S(u,v^{[q]})-\mathcal I_S(v,u^{[q]})\\
 &\geq a^d+b^d-ab^{d-1}-ba^{d-1}
       =(a-b)(a^{d-1}-b^{d-1})\geq|a-b|^d.
 \end{split}
 \tag{EV10}
\]
The last inequality follows by the binomial expansion after
ordering $a,b$. Thus (1) implies energy convergence and a
uniform energy bound. The capacity estimate (EV8) gives
all the asserted forward capacity conclusions.

Conversely suppose (2). In particular the energies are
uniformly bounded. For each bounded weight truncation and
each original compact cutoff, the proved local bounded-weight
capacity continuity gives
\[
 \begin{split}
 \liminf_j\mathcal I_S(u_j,u^{[q]})&\geq e_1(u),\\
 \liminf_j\mathcal I_S(u,u_j^{[q]})&\geq e_1(u).
 \end{split}
 \tag{EV11}
\]
Here is the order of limits: first replace the
undifferentiated weight by $-\max(\cdot,-L)$ and insert a
compact nonnegative cutoff; the original bounded-weight
continuity passes $j\to\infty$. Then let the cutoff
increase to one and $L\to\infty$, using positivity.
The derivative slots are controlled by the uniform power
energy hypothesis through the already proved top capacity
stability theorem. Thus neither line presupposes critical
weighted continuity.

The optimal Holder upper bounds are respectively
$e_1(u_j)^{1/d}e_1(u)^{q/d}$ and
$e_1(u)^{1/d}e_1(u_j)^{q/d}$, both tending to $e_1(u)$.
Thus both cross integrals converge to $e_1(u)$.
The first line of (EV10) now proves (1).
\end{proof}

\begin{theorem}[Critical mixed energy and weighted-current continuity]
\label{thm:r11-strong-critical-mixed-continuity}
Let $u_{i,j},u_i\in\mathcal E^1$ for $0\leq i\leq q$,
and assume $\mathscr I_S(u_{i,j},u_i)\to0$ for every slot.
Then
\[
 \mathcal I_S(u_{0,j},\ldots,u_{q,j})
                  \longrightarrow\mathcal I_S(u_0,\ldots,u_q),
 \tag{EV12}
\]
and the finite positive measures
\[
 (-u_{0,j})M(u_{1,j},\ldots,u_{q,j})
                      \longrightarrow(-u_0)M(u_1,\ldots,u_q)
 \tag{EV13}
\]
converge against every bounded original smooth scalar test
on $V$, including tests without compact support. In
particular there is neither a critical pole residue nor
escaped boundary energy. No common pointwise lower barrier
is required in this statement.
\end{theorem}

\noindent\textbf{Status: \texttt{NEW\_PROVED}}. The
strong discrepancy hypothesis is substantive; the critical
failure in (PS7) concerns mere capacity convergence with
a uniform energy bound.

\begin{proof}
By (EV9), every slot has uniformly bounded first energy.
Telescope the fully symmetric multilinear integral
$\mathcal I_S$ one slot at a time. Move each changing
slot to the undifferentiated position by (GE5).
The absolute value of that difference is bounded by
$\int|u_{i,j}-u_i|M(\mathbf z_j)$ for a mixed tuple
of unchanged or already changed slots. Its positive
expansion is at most $1/q!$ times the diagonal product
of their sum. That sum has uniformly bounded energy by
(GE11). Estimate (EV1) makes each of the finitely many
differences tend to zero. This proves (EV12).

Write $M_j=M(u_{1,j},\ldots,u_{q,j})$ and $M$ for its
limit. The same estimate gives
\[
                  \int|u_{0,j}-u_0|\,dM_j\longrightarrow0.
 \tag{EV14}
\]
For $u=u_0$, put $u^R=\max(u,-R)$. Since $u^R-u\geq0$,
the nonnegative discrepancy satisfies
\[
 0\leq\mathscr I_S(u,u^R)
       =\int(u^R-u)(d\mu_u-d\mu_{u^R})
       \leq\int(u^R-u)\,d\mu_u\longrightarrow0.
 \tag{EV15}
\]
The last limit is dominated convergence against the finite
first-energy weight $(-u)\mu_u$.
Apply (EV1) to $u,u^R$ and to the uniformly bounded-energy
sum of the differentiated slots of $M_j$. This proves
$\sup_j\int|u-u^R|\,dM_j\to0$; the same holds for $M$.
The constants are uniform because $e_1(u^R)\leq e_1(u)$
by energy monotonicity and all other energies are bounded.

For fixed $R$, local bounded-weight capacity continuity
gives $(-u^R)M_j\to(-u^R)M$ on every original compact
scalar test. Use the preceding uniform weight tail and
(EV14) to pass $R\to\infty$. This proves local convergence
in (EV13). Their total masses converge by (EV12).
Positive finite Radon measures with local convergence
and convergent full masses are tight: choose a compact
original cutoff whose integral in the limit is arbitrarily
close to the full limit mass, and pass its local integral.
This controls the complements uniformly for all large $j$.
It follows that every bounded original smooth scalar test,
with no compact-support restriction, has the same limit.
\end{proof}

\begin{corollary}[Energy-dual maximizing sequences are strongly Cauchy]
\label{cor:r11-dual-maximizing-Cauchy}
Let $\mu$ be a positive Radon datum with
$D=\mathfrak d_1(\mu)<\infty$ as defined in (DO1).
For $\mathcal F_\mu(u)=\mathscr A_S(u)-\int u\,d\mu$,
\[
 0\leq\sup_{u\in\mathcal E^1}\mathcal F_\mu(u)
               \leq\frac q{q+1}D^{(q+1)/q}<\infty.
 \tag{EV16}
\]
Every maximizing sequence is uniformly bounded in first
energy and is Cauchy for $\mathscr I_S$ and in every local
outer capacity. This does not assert the existence of an
actual classical limit or a classical Dirichlet solution.
\end{corollary}

\noindent\textbf{Status: \texttt{NEW\_PROVED}} as a
compactness estimate in the unbounded nonlinear class;
the representation and equation gates remain explicit.

\begin{proof}
Positive homogeneity and the definition of $D$ give
$\int(-u)\,d\mu\leq D e_1(u)^{1/d}$.
If $e_1(u)=0$, scale $u$ by an arbitrary positive number
in (DO1) to obtain the same zero conclusion.
Thus, with $x=e_1(u)^{1/d}$,
\[
                \mathcal F_\mu(u)\leq-x^d/d+Dx.
\]
Maximizing this real polynomial gives (EV16), and it tends
to $-\infty$ as $x\to\infty$. Since $0$ is admissible,
every tail of a maximizing sequence has
$\mathcal F_\mu(u_j)\geq-1$ and hence
\[
          e_1(u_j)\leq
             \max\{2d,(2dD)^{d/(d-1)}\}.
 \tag{EV17}
\]
Discarding finitely many terms gives uniform energy for
the sequence.

The frozen bounded common-collar midpoint argument now
applies to the proved global variation (GE16), giving
\[
 \mathscr A_S\left(\frac{u+v}2\right)
       -\frac{\mathscr A_S(u)+\mathscr A_S(v)}2
            \geq\frac q{4^{q+1}}\mathscr I_S(u,v).
 \tag{EV18}
\]
For completeness the midpoint Green formula is
$\frac q2\int_0^1\min(t,1-t)
 \int\Gamma_{N(u_t^{q-1})}(u-v)\,dt$.
On $[1/4,3/4]$, positive mixed expansion bounds its
energy integrand below by $4^{-(q-1)}\mathscr I_S(u,v)$.
The interval length and the scalar weight give exactly
the constant in (EV18). All global energy integrals
are finite by (GE7).

The datum integral is affine, so it cancels in the midpoint
gap. If $L=\sup\mathcal F_\mu$, then
\[
 \frac q{4^{q+1}}\mathscr I_S(u_j,u_l)
        \leq L-\frac{\mathcal F_\mu(u_j)+
                              \mathcal F_\mu(u_l)}2\longrightarrow0.
\]
The uniform energy bound and (EV8) give the all-local
outer-capacity Cauchy conclusion. Producing an actual
classical representative of that completed class, and
valid admissible variations that identify its curvature
with the datum, are not proved by this argument. The
compact-source obstruction (DO9) prevents omitting those
gates in a universal classical existence assertion.
\end{proof}

\begin{corollary}[Strong first-energy convergence controls every graded endpoint]
\label{cor:r11-strong-first-energy-graded-endpoints}
Let $u_{i,j},u_i\in\mathcal E^1$ be a finite family with
$\mathscr I_S(u_{i,j},u_i)\to0$ in every slot. For
$w_j=\sum_i u_{i,j}$ one has
\[
 \lim_{R\to\infty}\sup_j\int_R^\infty
                 t^q C_S(\{w_j<-t\};V)\,dt=0.
 \tag{EV19}
\]
Consequently all three convergence assertions (GS3), and
the double-prime first-order assertion, hold with
$\sigma=1$ and weak inequalities in their degree/weight
conditions. They hold on every original compact
complementary test and every actual moment refinement,
without a common pointwise lower barrier. These are
original-current convergence assertions; total-variation
convergence for arbitrary varying tuples is not claimed.
\end{corollary}

\noindent\textbf{Status: \texttt{NEW\_PROVED}} as the
uniform-tail consequence of the strong topology; applying
the established endpoint theorem is a direct consequence.

\begin{proof}
First consider one sequence $u_j\to u$ in discrepancy.
Put $a_j=-u_j$ and $b=-u$ off the joint polar set.
The uniform energy bound follows from (EV9). On
$\{a_j>R\}$ split according as $b\leq a_j/2$ or
$b>a_j/2$. In the first region $a_j\leq2|u_j-u|$;
in the second, $a_j<2b$ and $b>R/2$.
For $b>R/2$ one has
$b\leq2(b-R/4)^+=2|u-u^{R/4}|$.
Thus
\[
 \int_{\{a_j>R\}}a_j\,d\mu_{u_j}
    \leq2\int|u_j-u|\,d\mu_{u_j}
             +4\int|u-u^{R/4}|\,d\mu_{u_j}.
 \tag{EV20}
\]
By (EV1), the first term tends to zero as $j\to\infty$.
By (EV1) and (EV15), the second tends to zero as
$R\to\infty$ uniformly in $j$. To obtain the supremum
over all $j$, first make the first term small for
$j\geq j_0$; the finitely many earlier rows each have
integrable first-energy weight, so their tails tend
to zero separately. Then let $R\to\infty$, and finally
$j_0\to\infty$. This proves uniform vanishing of the
left side of (EV20).

The sublevel capacity bound (PE4), before its full
integration, gives for every row
\[
 \int_R^\infty t^q C_S(\{u_j<-t\};V)\,dt
       \leq2^{q+1}\int(a_j-R/2)^+\,d\mu_{u_j}.
 \tag{EV21}
\]
Its right side is bounded by a constant times the
uniformly vanishing energy tail at $R/2$.
This proves (EV19) for one row.

For $m$ slots, the actual sublevel inclusion
$\{w_j<-t\}\subset\bigcup_{i=1}^m\{u_{i,j}<-t/m\}$
and subadditivity imply that its capacity-tail integral
is at most $m^{q+1}$ times the sum of the individual
tail integrals starting at $R/m$. This proves (EV19)
for the family. The full local capacity convergence is
already supplied by (EV8), so every hypothesis of
Theorem~\ref{thm:r11-nonmonotone-graded-stability}, including
its endpoint condition (GS4), now holds at $\sigma=1$.
Applying that theorem proves the current conclusions.
\end{proof}
